\subsection{Plurality Method}

\content{% main
\begin{definition} A \textbf{preference ballot} is a ballot in which each voter ranks the
choices according to preference. 
\end{definition}
}{% presentation
\vspace{2in}
}{% nopresentation
\vspace{2in}
}{% comments
Draw out what a simple preference ballot for one voter might look like. What
does the total data look like? Note that we usually combine the preference
ballots that look similar. 
}

\content{
\begin{definition} A \textbf{Preference schedule} is a condensed way of
writing out several preference ballots. 
\end{definition}
}{% presentation
\vspace{2in}
}{% nopresentation
\vspace{2in}
}{% comments
Draw an example. Notice that the total number of votes is the sum of the numbers
above each column. 
}

\content{
        \begin{example} Let's look the following vote. \\
        \begin{tabular}{|l|c|c|c|c|}
        \hline
        & 1 & 3 & 3 & 3 \\
        \hline
        1st Choice  & A & A & O & H \\
        \hline
        2nd Choice  & O & H &H&A \\
        \hline
        3rd Choice & H&O&A&O \\
        \hline
        \end{tabular}
        \end{example}
}{% presentation
\vspace{3in}
}{% nopresentation
\vspace{2in}
}{% comments
How do we determine who wins? \\
Introduce the plurality method!\\
How many voters are there? Who has the most 1st
place votes?  
}


\content{
        \begin{example} Now you try it!\\
        \begin{tabular}{|l|c|c|c|c|}
        \hline
        & 10& 8 & 4 & 3 \\
        \hline
        1st choice & A & B & C & B \\
        \hline
        2nd choice & B & C & B & A \\
        \hline
        3nd choice & C & A & A & C \\
        \hline
        \end{tabular}
        \begin{enumerate}
        \item How many voters are there? 
        \item How many first place votes does $B$ have? 
        \item Who wins under the plurality method? 
        \end{enumerate}
        \end{example}
}{% presentation
\vspace{1in}
}{% nopresentation
\vspace{1in}
}{% comments
}

\content{
        \begin{example} Let's look the vote from before. \\
        \begin{tabular}{|l|c|c|c|c|}
        \hline
        & 1 & 3 & 3 & 3 \\
        \hline
        1st Choice  & A & A & O & H \\
        \hline
        2nd Choice  & O & H &H&A \\
        \hline
        3rd Choice & H&O&A&O \\
        \hline
        \end{tabular}
        \end{example}
        What could go wrong here? 
}{% presentation
\vspace{4in}
}{% nopresentation
\vspace{2in}
}{% comments
What happens if one of the options is no longer available? \\
Start to do a pairwise comparison. Notice that when $H$ is compared to $A$, $H$
seems to be preffered, but did not win. Is that fair? What does it mean to have
a fair election?
}


\content{
\begin{definition} A \textbf{fairness criterion} is a statement which seems like
it should be true during an election.
\end{definition}
}{% presentation
}{% nopresentation
}{% comments
}

\content{
We have now run into our very first Fairness Criterion!\\

\begin{definition} \textbf{(Condercet Criterion)} If there is a choice which is
preffered in every one-to-one comparison to the other choices, that choice
should win the election. That choice (if it exists) is called the
\textbf{Condocet Candidate}.
\end{definition}
}{% presentation
}{% nopresentation
}{% comments
}


\content{
\begin{example} Your turn again! Which option wins the following vote using the
plurality method? Is there a Condorcet Candidate? If there is one, did they win
the election? \\ 
\begin{tabular}{|l|c|c|c|c|c|}
\hline
& 44 & 14 & 70 & 22 & 80 \\
\hline
1st Choice & G & G &M&M&B\\
\hline
2nd Choice &M&B&G&B&M\\
\hline
3rd Choice &B&M&B&G&G\\
\hline
\end{tabular}
\end{example}
}{% presentation
\vspace{3in}
}{% nopresentation
\vspace{2in}
}{% comments
}
